\scanpage{176}
\mspage{125}
% archive sheet number: 176
canonique du diagramme \unsure{($*$)|($\alpha$)} dans ($**$), qui est l'identité \unsure{sur|pour} les sommets $X_0$ et $\unsure{\mathscr{X}|\mathcal X}$, et qui \deleted{provient du fait que} \unsure{sur|pour} $X_1,X_2$ est donné par les inclusions canoniques

\[
\begin{aligned}
X_1&\longrightarrow{}_1\backslash\mathscr{X}_{\deleted{/1}}\deleted{\unsure{{}_{/1}}},\qquad&
X_2&\longrightarrow{}_2\backslash\mathscr{X}_{\deleted{/2}}.\\
\Vert&&\Vert\\
\mathscr{X}_1\ \text{fibre}&&\mathscr{X}_2\ \text{fibre}\\
\text{de }\mathscr{X}\text{ en }1&&\text{de }\mathscr{X}\text{ en }2
\end{aligned}
\]

Ces inclusions, on \deleted{\unsure{le}} le sait, sont $\unsure{W|W_0}$-asphériques (elles ont $=$ un adjoint à droite), à fortiori elles sont dans $W$. Donc le diagramme non commutatif ($*$) est isomorphe, en tant que diagramme dans $\mathrm{Hot}_W$, au diagramme commutatif cartésien et cocartésien \inserted{($**$)} formé d'immersions \deleted{\unsure{ouvertes}}\inserted{fermées}. \deleted{\unsure{C'est|Ce}}Comme $\mathscr{X}$ est bel et bien somme amalgamée de ses deux \deleted{ouverts}\inserted{fermés} $\,{}_1\backslash\mathscr{X}_{\deleted{/1}},{}_2\backslash\mathscr{X}_{\deleted{/2}}$ sur leur intersection $X_0$, cela étaye l'intuition qui \deleted{\unsure{fait}} voit $\mathscr{X}$ comme une \lq\lq somme amalgamée \unsure{\lq\lq}homotopique\rq\rq. \unsure{$\S$}

Donnons \deleted{tout} de suite l'application suivante :

\uline{Proposition 1\deleted{\unsure{2|3}}\unsure{3|2}} : Soit un diagramme

\diagram{\begin{tikzcd}[ampersand replacement=\&,row sep=0pt,column sep=small]
\& X_1\\
X_0\arrow[ur,"\alpha_1"]\arrow[dr,"\alpha_2"'] \&\\
\& X_2
\end{tikzcd}}
dans $\mathrm{Cat}$, d'où le diagramme \unsure{($\alpha$)|($*$)} (p.~123)

\diagram{\begin{tikzcd}[ampersand replacement=\&,row sep=0pt,column sep=small]
\& X_1\arrow[dr,"i_1"] \&\\
X_0\arrow[ur,"\alpha_1"]\arrow[dr,"\alpha_2"'] \& \& \mathscr{X}\\
\& X_2\arrow[ur,"i_2"'] \&
\end{tikzcd}}
